\chapter{Code Generation}

\section{Overview}

The code generator traverses the annotated parse tree of a valid WLP4 program and emits MIPS
assembly. There are infinitely many correct translations; our goal is a \emph{simple,
correct} one, structured so each construct's code can be generated independently
(syntax-directed translation). We write $\mathrm{code}(X)$ for the assembly produced for
subtree $X$.

\begin{important}[Invariant]
$\mathrm{code}(\mathit{expr})$ leaves the value of \emph{expr} in \code{\$3}, and leaves every
other register (except scratch \code{\$5}--\code{\$7}) and the stack as they were.
\end{important}

\section{Conventions}

\begin{center}
\begin{tabular}{ll}
\toprule
register & use \\
\midrule
\code{\$0} & zero \\
\code{\$1}, \code{\$2} & parameters of \code{wain} (and arguments to runtime routines) \\
\code{\$3} & result of every expression and procedure \\
\code{\$4} & constant 4 \\
\code{\$5}--\code{\$7} & scratch \\
\code{\$11} & constant 1 \\
\code{\$29} & frame pointer \\
\code{\$30} & stack pointer \\
\code{\$31} & return address \\
\bottomrule
\end{tabular}
\end{center}

Helper ``macros'' the generator emits often:
\begin{lstlisting}[style=mips]
; push($r)
sw $r, -4($30)
sub $30, $30, $4
; pop($r)
add $30, $30, $4
lw $r, -4($30)
\end{lstlisting}

\section{Variables and the stack frame}

All variables live on the stack, addressed relative to a \textbf{frame pointer} \code{\$29}
that is fixed for the whole procedure (unlike \code{\$30}, which moves as temporaries are
pushed). The symbol table records each variable's offset.

\paragraph{Prologue of wain.} Import the runtime, set up constants, and push the parameters and locals:
\begin{lstlisting}[style=mips]
.import print
.import init
.import new
.import delete
lis $4
.word 4
lis $11
.word 1
sub $29, $30, $4     ; $29 points at first variable slot
sw $1, -4($30)       ; a at offset 0
sub $30, $30, $4
sw $2, -4($30)       ; b at offset -4
sub $30, $30, $4
; locals: int x = 7;  -> offset -8
lis $5
.word 7
sw $5, -4($30)
sub $30, $30, $4
\end{lstlisting}
The $k$-th variable (counting from 0) is at offset $-4k$ from \code{\$29}.
Reading a variable: \code{lw \$3, off(\$29)}; writing: \code{sw \$3, off(\$29)}.
\begin{center}
\begin{tikzpicture}[>=Stealth,cell/.style={draw,minimum width=3.2cm,minimum height=0.55cm,font=\small}]
  \node[cell,fill=gray!10] (t) at (0,0) {loader's stack};
  \node[cell,below=0pt of t,fill=blue!5] (a) {\code{a} (from \code{\$1})};
  \node[cell,below=0pt of a,fill=blue!5] (b) {\code{b} (from \code{\$2})};
  \node[cell,below=0pt of b,fill=blue!5] (x) {\code{x} $=7$};
  \node[cell,below=0pt of x,fill=gray!5,draw=gray] (f) {\textcolor{gray}{free}};
  \foreach \c/\o in {a/0,b/-4,x/-8} \node[font=\scriptsize,anchor=east] at ([xshift=-2pt]\c.west) {offset $\o$};
  \draw[<-] (a.east) -- ++(0.8,0) node[right]{\code{\$29} (fixed)};
  \draw[<-] (x.east) -- ++(0.8,0) node[right]{\code{\$30} (moves)};
\end{tikzpicture}

\small \code{wain}'s frame after the prologue (higher addresses at the top).
\end{center}
\paragraph{Epilogue.} \code{add \$30, \$29, \$4} (pop the frame) then \code{jr \$31}.

\section{Expressions}

\subsection{Constants and variables}
\begin{lstlisting}[style=mips]
; code(NUM)              ; code(ID)
lis $3                   lw $3, off($29)
.word NUM
\end{lstlisting}
(\code{NULL} is represented by the value 1, not 0: it is not word-aligned, so dereferencing
NULL crashes the emulator, which is the desired behaviour. Use \code{add \$3, \$0, \$11}.)

\subsection{Binary operators}
Use the stack to hold the left operand while computing the right one:
\begin{lstlisting}[style=mips]
; code(e1 + e2), both int
code(e1)
push($3)
code(e2)
pop($5)
add $3, $5, $3
\end{lstlisting}
Similarly \code{sub}; for \code{*} use \code{mult \$5, \$3; mflo \$3}; for \code{/} use
\code{div \$5, \$3; mflo \$3}; for \code{\%} use \code{mfhi \$3}. This recursive scheme works
for arbitrarily nested expressions because each subexpression restores the stack.

\begin{aside}
An easy optimisation: if \code{e2} is a constant or variable, compute it straight into a
register without the push/pop. Another: allocate a different temporary register per nesting
level and only spill to the stack when they run out.
\end{aside}

\subsection{Pointer arithmetic}
Types (from semantic analysis) decide the code:
\begin{center}
\small
\begin{tabular}{ll}
\toprule
expression & code after \code{\$5} $=e_1$, \code{\$3} $= e_2$ \\
\midrule
\code{int* + int} & \code{mult \$3, \$4; mflo \$3; add \$3, \$5, \$3} \\
\code{int + int*} & \code{mult \$5, \$4; mflo \$5; add \$3, \$5, \$3} \\
\code{int* - int} & \code{mult \$3, \$4; mflo \$3; sub \$3, \$5, \$3} \\
\code{int* - int*} & \code{sub \$3, \$5, \$3; div \$3, \$4; mflo \$3} \\
\bottomrule
\end{tabular}
\end{center}

\subsection{Pointers: \code{*}, \code{\&}, and assignment through pointers}
\begin{lstlisting}[style=mips]
; code(*e)           ; code(&ID)              ; code(&(*e)) = code(e)
code(e)              lis $3
lw $3, 0($3)         .word off
                     add $3, $3, $29
\end{lstlisting}
Assignment:
\begin{lstlisting}[style=mips]
; ID = e;            ; *e1 = e2;
code(e)              code(e1)          ; address
sw $3, off($29)      push($3)
                     code(e2)          ; value
                     pop($5)
                     sw $3, 0($5)
\end{lstlisting}

\section{Tests}

Produce 1 (true) or 0 (false) in \code{\$3}. Use \code{slt} for \code{int} and
\code{sltu} for pointers (addresses are unsigned).
\begin{center}
\small
\begin{tabular}{ll}
\toprule
test & code after \code{\$5} $=e_1$, \code{\$3} $=e_2$ \\
\midrule
$e_1 < e_2$ & \code{slt \$3, \$5, \$3} \\
$e_1 > e_2$ & \code{slt \$3, \$3, \$5} \\
$e_1 \ne e_2$ & \code{slt \$6, \$3, \$5; slt \$7, \$5, \$3; add \$3, \$6, \$7} \\
$e_1 = e_2$ & ($\ne$ code), then \code{sub \$3, \$11, \$3} \\
$e_1 \ge e_2$ & \code{slt \$3, \$5, \$3; sub \$3, \$11, \$3} \\
$e_1 \le e_2$ & \code{slt \$3, \$3, \$5; sub \$3, \$11, \$3} \\
\bottomrule
\end{tabular}
\end{center}

\section{Control flow}

Each \code{if} and \code{while} needs unique labels; keep a global counter.
\begin{lstlisting}[style=mips]
; if (T) {S1} else {S2}       ; while (T) {S}
code(T)                       loop7:
beq $3, $0, else4             code(T)
code(S1)                      beq $3, $0, endWhile7
beq $0, $0, endif4            code(S)
else4:                        beq $0, $0, loop7
code(S2)                      endWhile7:
endif4:
\end{lstlisting}
A sequence of statements is just the concatenation of their code.

\subsection{Printing}
\code{print} (from \code{print.merl}) prints \code{\$1} as a decimal number and a newline.
Since \code{\$1} holds \code{wain}'s parameter we save it:
\begin{lstlisting}[style=mips]
; println(e);
push($1)
code(e)
add $1, $3, $0
push($31)
lis $5
.word print
jalr $5
pop($31)
pop($1)
\end{lstlisting}

\section{Procedures}

\paragraph{Naming.} Prefix every WLP4 procedure label, e.g.\ \code{Ffoo}, so it cannot
clash with generated labels (\code{else4}) or runtime routines (\code{print}, \code{new}).
\code{wain}'s code comes first in the output so execution starts there; other procedures
follow.

\paragraph{Caller.} Saves \code{\$29} and \code{\$31}, pushes arguments in order, calls,
pops the arguments, restores.
\begin{lstlisting}[style=mips]
; code(f(e1, ..., en))
push($29)
push($31)
code(e1)
push($3)
...
code(en)
push($3)
lis $5
.word Ff
jalr $5
; pop n arguments (n adds, or one add of 4n)
pop($31)
pop($29)
\end{lstlisting}

\paragraph{Callee.} Sets up its frame, runs its body, and tears the frame down:
\begin{lstlisting}[style=mips]
Ff:
sub $29, $30, $4     ; new frame pointer
; push locals (initial values)
; push registers this procedure will modify ($5, $6, $7, ...)
code(statements)
code(return expr)    ; result in $3
; pop saved registers
add $30, $29, $4     ; discard locals
jr $31
\end{lstlisting}

\paragraph{Offsets.} With $n$ parameters, the caller pushed $a_1,\dots,a_n$, so after
\code{sub \$29, \$30, \$4}, parameter $a_i$ is at offset $4(n-i+1)$ (positive) from
\code{\$29}: $a_n$ at $4$, $a_1$ at $4n$. Locals are at $0, -4, -8,\dots$ as in \code{wain}.
The symbol table is filled with these offsets while processing the procedure's declarations.

\begin{center}
\begin{tikzpicture}[cell/.style={draw,minimum width=3.6cm,minimum height=0.55cm,font=\small}]
  \node[cell] (a) at (0,0) {saved \code{\$29}};
  \node[cell,below=0pt of a] (b) {saved \code{\$31}};
  \node[cell,below=0pt of b] (c) {$a_1$ \quad (off $4n$)};
  \node[cell,below=0pt of c] (d) {$\vdots$};
  \node[cell,below=0pt of d] (e) {$a_n$ \quad (off 4)};
  \node[cell,below=0pt of e,fill=blue!5] (f) {local 1 \quad (off 0)};
  \node[cell,below=0pt of f,fill=blue!5] (g) {local 2 \quad (off $-4$)};
  \node[cell,below=0pt of g,fill=gray!10] (h) {saved registers};
  \draw[Stealth-] (f.east) -- ++(1.2,0) node[right]{\code{\$29}};
  \draw[Stealth-] (h.west) -- ++(-1.2,0) node[left]{\code{\$30}};
  \draw ($(a.north east)+(0.2,0)$) -- node[right,font=\small]{caller pushes} ($(e.south east)+(0.2,0)$);
  \draw ($(f.north west)+(-0.2,0)$) -- node[left,font=\small]{callee pushes} ($(h.south west)+(-0.2,0)$);
\end{tikzpicture}
\end{center}

\begin{important}[Who saves what]
Here the caller saves \code{\$29} and \code{\$31}; the callee saves any other register it
modifies. \code{\$4}, \code{\$11} are never modified. Mixing up the order of pushes and pops
is the most common bug: draw the stack!
\end{important}

\section{Dynamic memory: \code{new} and \code{delete}}

The runtime \code{alloc.merl} provides \code{init}, \code{new} and \code{delete}.
\begin{itemize}
  \item \code{init} must be called once in \code{wain}'s prologue (with \code{\$31} saved). It
        expects \code{\$1}/\code{\$2} as given by \code{mips.array}; if \code{wain}'s first
        parameter is \code{int} (\code{mips.twoints}), set \code{\$2} to 0 first.
  \item \code{new}: number of words in \code{\$1}; returns address in \code{\$3}, or 0 on failure.
        WLP4 must return NULL on failure, so convert 0 to 1.
  \item \code{delete}: address in \code{\$1}. Skip the call when the pointer is NULL.
\end{itemize}

\begin{lstlisting}[style=mips]
; code(new int[e])
code(e)
add $1, $3, $0
push($31)
lis $5
.word new
jalr $5
pop($31)
bne $3, $0, 1      ; success: skip next instruction
add $3, $11, $0    ; failure: return NULL (1)
\end{lstlisting}
\begin{lstlisting}[style=mips]
; code(delete [] e;)
code(e)
beq $3, $11, skipDelete9   ; never free NULL
add $1, $3, $0
push($31)
lis $5
.word delete
jalr $5
pop($31)
skipDelete9:
\end{lstlisting}
(Also save \code{\$1} around these if it is still needed.)

\section{Putting it together}

The output is assembled to MERL and linked with the runtime:
\begin{lstlisting}[style=bash]
wlp4scan < prog.wlp4 | wlp4parse | wlp4gen > prog.asm
cs241.linkasm < prog.asm > prog.merl
cs241.linker prog.merl print.merl alloc.merl > exec.mips
mips.twoints exec.mips      # or mips.array
\end{lstlisting}

\begin{example}[Complete output for a small program]
Imports and the call to \code{init} are omitted since this program uses no runtime routines.
\begin{lstlisting}[style=cpp]
int wain(int a, int b) { return a + b * 2; }
\end{lstlisting}
\begin{lstlisting}[style=mips]
lis $4
.word 4
lis $11
.word 1
sub $29, $30, $4
sw $1, -4($30)
sub $30, $30, $4     ; a at 0
sw $2, -4($30)
sub $30, $30, $4     ; b at -4
lw $3, 0($29)        ; a
sw $3, -4($30)
sub $30, $30, $4     ; push
lw $3, -4($29)       ; b
sw $3, -4($30)
sub $30, $30, $4     ; push
lis $3
.word 2
add $30, $30, $4
lw $5, -4($30)       ; pop b
mult $5, $3
mflo $3              ; b*2
add $30, $30, $4
lw $5, -4($30)       ; pop a
add $3, $5, $3       ; a + b*2
add $30, $29, $4
jr $31
\end{lstlisting}
\end{example}

How the stack evolves while the body computes \code{a + b * 2}:
\begin{center}
\begin{tikzpicture}[>=Stealth,cell/.style={draw,minimum width=2cm,minimum height=0.5cm,font=\small},
  tmp/.style={cell,fill=orange!12}]
  \foreach \x/\i/\t in {0/0/{after \code{push(a)}},3.4/1/{after \code{push(b)}},6.8/2/{after both pops}} {
    \node[font=\scriptsize] at (\x,0.55) {\t};
    \node[cell,fill=blue!5] (a\i) at (\x,0) {\code{a}};
    \node[cell,fill=blue!5,below=0pt of a\i] (b\i) {\code{b}};
  }
  \node[tmp,below=0pt of b0] (p0) {copy of \code{a}};
  \node[tmp,below=0pt of b1] (p1) {copy of \code{a}};
  \node[tmp,below=0pt of p1] (q1) {copy of \code{b}};
  \draw[<-] (p0.west) -- ++(-0.5,0) node[left,font=\scriptsize]{\code{\$30}};
  \draw[<-] (q1.west) -- ++(-0.5,0) node[left,font=\scriptsize]{\code{\$30}};
  \draw[<-] (b2.west) -- ++(-0.5,0) node[left,font=\scriptsize]{\code{\$30}};
  \node[font=\scriptsize,below=4pt of p0] {\code{\$3} $=$ \code{a}};
  \node[font=\scriptsize,below=4pt of q1] {\code{\$3} $=2$, then \code{b*2}};
  \node[font=\scriptsize,below=4pt of b2] {\code{\$3} $=$ \code{a + b*2}};
\end{tikzpicture}
\end{center}
Each binary operator pushes exactly one temporary and pops it again, so the stack is back to
the bare frame when the expression is done.

\section{Summary}
\begin{itemize}
  \item Syntax-directed translation: each construct's code is built from its children's code.
  \item Every expression leaves its result in \code{\$3}; temporaries go on the stack.
  \item Variables are at fixed offsets from the frame pointer \code{\$29}.
  \item Types select instructions: scale by 4 for pointer arithmetic, \code{sltu} for pointer comparisons.
  \item Unique labels for \code{if}/\code{while}; prefix procedure names.
  \item Calls: caller saves \code{\$29}, \code{\$31} and pushes arguments; callee sets
        \code{\$29}, pushes locals and saved registers.
  \item \code{new}/\code{delete}/\code{print} come from runtime MERL modules linked in; NULL is 1.
\end{itemize}
